% ============================================================
% Single-column article version
% "Accessible Brain Morphometry from Point-of-Care
%  Ultra-Low-Field MRI Using Physics-Constrained Deep Learning"
% ============================================================
\documentclass[12pt,a4paper]{article}
\usepackage[margin=2.5cm]{geometry}
\usepackage{cite}
\usepackage{amsmath,amssymb,amsfonts}
\usepackage{graphicx}
\usepackage{textcomp}
\usepackage{booktabs}
\usepackage{multirow}
\usepackage{xcolor}
\usepackage{hyperref}
\usepackage{setspace}
\usepackage{titlesec}
\graphicspath{{paper_figures/}}

\onehalfspacing

\begin{document}

\title{\textbf{Accessible Brain Morphometry from Point-of-Care\\
Ultra-Low-Field MRI Using Physics-Constrained Deep Learning}}

\author{Author Name\\
\small Department Name, Institution Name, City, Country\\
\small \texttt{author@institution.edu}}

\date{}
\maketitle

\noindent\textbf{Funding:} Funding acknowledgment placeholder.
The OASIS-1 dataset is available at \url{https://www.oasis-brains.org}.
The ds006557 dataset is available at \url{https://openneuro.org/datasets/ds006557}.
Code will be released upon acceptance.

\bigskip

% ── Abstract ─────────────────────────────────────────────────────────────────
\begin{abstract}
Normalized whole-brain volume (nWBV) is an extensively studied biomarker for
dementia progression but requires costly, high field MRIs, along with compute
intensive segmentation pipelines to obtain the same routinely. We propose
ViT3D (Vision Transformer 3D, 4.23M params), a small architecture capable of
directly predicting nWBV from a single axial T2 weighted acquisition using a
64\,mT Hyperfine Swoop bed-side MRI scanner -- which can be bought for about
10$\times$ less cost compared to 3\,T MRI. The pretraining is done using
physics constrained simulations of 64\,mT image formation process (tissue
specific T1/T2 relaxations, Rician noise, B0 inhomogeneity), done on 156 IXI
participants, followed by finetuning on 375 OASIS-1 participants using their
nWBV annotations generated using FreeSurfer. On OASIS-1 held-out test data
($n = 38$) Pearson $r$ is 0.722 [95\% CI: 0.524--0.847, $p < 0.001$] for
this model, and MAE is 0.058. It is able to distinguish subjects based on
severity of their dementia diagnosis (CDR 0 $\to$ 0.752, CDR 0.5 $\to$
0.739, CDR 1.0 $\to$ 0.678 mean nWBV). On a separate test set of 23 actual
Hyperfine 64\,mT MRI scans (ds006557), the model predicts age-related brain
volume reduction (Spearman $\rho = -0.597$, $p = 0.003$) and reaches a MAE
of 0.010\,nWBV following lightweight domain adaptation -- below the accepted
clinical threshold of 0.020. An ablation study proves that physics-guided
simulation exceeds image blurring ($r = 0.949$ vs.\ $0.931$, $\Delta r =
+0.018$) under equivalent training scenarios. Inference takes 47\,ms and
does not require FreeSurfer installation, enabling edge deployment.
\end{abstract}

\noindent\textbf{Keywords:}
Brain morphometry, dementia biomarker, low-field MRI, normalized whole-brain
volume, physics-informed simulation, point-of-care neuroimaging,
Vision Transformer

% ============================================================
\section{Introduction}
\label{sec:introduction}
% ============================================================

Dementia afflicts more than 55 million individuals globally,
with cases expected to rise three-fold by 2050~\cite{who2023}. Normalized
whole-brain volume (nWBV), which is considered the fraction of total brain
parenchymal tissue against total estimated intracranial volume, serves as a
highly validated measure of neurodegeneration, where every CDR stage accounts
for about 5--10\% drop in the nWBV~\cite{fotenos2005,buckner2004}. Monitoring
of nWBV allows identification of neurodegenerative processes prior to symptom
onset, hence it provides an excellent outcome measure for screening studies
and clinical trials~\cite{jack2018nia}.

nWBV measurement in accordance with the standardized protocol entails the use
of high-field (1.5\,T and/or 3\,T) MRI coupled with automated software tools
like FreeSurfer~\cite{fischl2012freesurfer} and FSL~\cite{smith2002fsl}.
However, these processes require computation time of 20--90 minutes per
patient, software packages larger than 12\,GB, and MRI systems priced at
USD~1--3\,million. Unfortunately, such advanced infrastructure is out of reach
for LMICs, rural communities, and point-of-care environments~\cite{prince2016world}.

The ultra-low field (ULF) MRI scanner like the Hyperfine Swoop (64\,mT) is
transportable, does not require shielding, and costs around \$50,000 without
any special room specifications~\cite{hyperfine2021}. It has already been used
in bedside stroke triage~\cite{kimberly2023bedside}, but its potential use in
quantifying brain morphology remains unexplored. The fundamental problem is
that 64\,mT scans have very different signal properties (contrast, resolution
of $1.6\times1.6\times5.0$\,mm, noise type---Rician noise, poor B0 shimming)
compared to 3\,T scans for which all current morphometry techniques have been
designed.

Previous methods involve: super-resolution upscaling of the ULF images
followed by use of existing pipelines~\cite{iglesias2023synthsr} with an
additional delay of 30--120\,seconds; domain
adaptation~\cite{chen2019synergistic}, which demands huge sets of matched
data, which is not possible for ULF field strengths; and SynthSeg+ for
contrast-agnostic segmentations~\cite{billot2023synthseg}, which is a
generalized model but takes 2--3\,minutes to process.

\textbf{Contributions.} We introduce, to our knowledge:
\begin{enumerate}
  \item This is among the first works applying a 3D
        Vision Transformer for direct nWBV regression from 64\,mT MRI,
        achieving inference in under 47\,ms without any additional software
        dependencies.
  \item A magnetic physics-informed 64\,mT simulator based on actual scanner
        parameters (T1/T2 of tissue at 64\,mT, Rician noise, and B0
        inhomogeneity) demonstrated better performance than Gaussian blurring
        degradation with $\Delta r = +0.018$.
  \item Real-world validation on the Hyperfine Swoop system against
        ground-truth data provided by FastSurfer ($n=23$), which showed
        clinically adequate mean absolute error (MAE: 0.010) with light
        domain adaptation.
  \item Failure analysis revealing atrophic brains and older patients as the
        main sources of failure, which is helpful in guiding future
        developments.
\end{enumerate}

% ============================================================
\section{Related Work}
\label{sec:related}
% ============================================================

\subsection{Brain Morphometry from MRI}

FreeSurfer~\cite{fischl2012freesurfer} and
FastSurfer~\cite{henschel2020fastsurfer} represent two leading techniques for
brain volumetric morphometry that rely on high-field acquisitions of isotropic
sub-millimeter T1-weighted volumes. There have been some attempts to perform
regression of volumetric biomarkers directly from MRI scans without using any
segmentation step. Peng et al.~\cite{peng2021brain} showed an effective
regression of brain age using structural MRI and avoiding segmentation
altogether, thus delivering comparable performance to the FreeSurfer based
reference ages. Roy et al.~\cite{roy2019quicknat} introduced QuickNAT -- a
fully convolutional network for fast brain segmentation -- while still
resorting to a pre-processing stage of brain segmentation before estimating
the brain volume. In the case of brain age regression, one could use 3D CNNs
on 3\,T data that deliver results within a 3--5 year mean absolute error
range~\cite{cole2017predicting,jonsson2019brain}. Rebsamen et
al.~\cite{rebsamen2020cortical} proved successful cortical thickness estimation
via deep learning segmentation, thus demonstrating competitiveness of
end-to-end regression techniques against pipeline-based techniques. To the
best of our knowledge, there is no previous research that has achieved a
direct regression of nWBV from ultra-low-field MRI.

\subsection{Ultra-Low-Field MRI}

V{\'a}{\v s}a et al.~\cite{vasa2025imaneur,vasa2025} describe test-retest
reliability and high-field MRI correlation for ULF brain morphometry and make
the ds006557 publicly available dataset ($n = 23$). Islam et
al.~\cite{islam2025aibrain} prove that deep neural networks (SynthSR,
HiLoResGAN) can minimize volumetric bias between 64\,mT and 3\,T
measurements. Arnold et al.~\cite{arnold2022simulated} report
simulation-based evaluation of diagnostic performance at 64\,mT using
publicly available data, which supports our simulation pre-training strategy.
Kimberly et al.~\cite{kimberly2023bedside} and Sheth et
al.~\cite{sheth2021assessmentultralow} validate the use of ULF MRI for
bedside neurology. Javadi et al.~\cite{javadi2025diffusion} compare diffusion
models and CNNs for 64\,mT-to-3\,T FLAIR synthesis, whereas Dayarathna et
al.~\cite{dayarathna2024adversarial} leverage an adversarial diffusion model
for ULF-to-HF translation. Gopinath et al.~\cite{gopinath2025cortical} prove
robust cortical mapping from low-field MRI, confirming the potential of
morphometry in ULF MRI. Baljer et al.~\cite{baljer2024paediatric} propose a
multi-orientation super-resolution U-Net for pediatric ULF MRI in low and
middle-income countries, proving the capability of deep learning to extract
structural information from bedside 64\,mT MRI. Minocha and
Yang~\cite{minocha2024lfmri} introduce a deep learning pipeline for early
Alzheimer disease diagnosis using LF MRI. None of the previous studies
quantify nWBV from 64\,mT scans without segmentation.

SynthSeg+~\cite{billot2023synthseg} is currently the best algorithm for
segmentation without considering contrast, achieving $r = 0.918$
compared to FreeSurfer nWBV on actual Hyperfine imaging data (this study). It
is used as the reference rather than the competitor since it takes 2--3
minutes and needs the entire installation of FreeSurfer, which cannot be
colocated with bedside equipment.

\subsection{Vision Transformers in Medical Imaging}

Dosovitskiy et al.~\cite{dosovitskiy2020image} proved that Vision Transformers
perform just as well or better than CNNs in terms of recognizing objects in
natural images. In the realm of medical images,
TransUNet~\cite{chen2021transunet} and Swin
UNETR~\cite{hatamizadeh2021swin} showed how Vision Transformers have an edge
over Convolutional Neural Networks in 3D image segmentation. With regard to
global regression tasks such as nWBV, which requires understanding the ratio
of brain tissues throughout the whole volume, the Vision Transformer has the
right architecture for the job since it uses global self-attention.

\subsection{Physics-Informed MRI Simulation}

Cross-domain MRI pretraining with simulated data usually involves Gaussian
blur as an approximation of degradation of
resolution~\cite{zhao2021superresolution}. Domain adaptation as an issue of
variability of scanners and acquisition protocols is well-known for brain MRI
studies; according to Karani et al.~\cite{karani2018lifelong}, models trained
using one scanner perform poorly on other scanners, and harmonization is
necessary as concluded by Dewey et al.~\cite{dewey2019deepharmony} for
accurate cross-site volumetry. This is one reason why we introduced a
lightweight domain adaptation module (Section~\ref{sec:training}).
SynthSR~\cite{iglesias2023synthsr} adopts a principled approach of
synthesizing images for super-resolution. In this study, our simulation
involves modeling the tissue contrast for 64\,mT using T1/T2 relaxation,
Rician noise, and inhomogeneous B0 field based on previous Hyperfine Swoop
measurements~\cite{vasa2025}.

% ============================================================
\section{Methods}
\label{sec:methods}
% ============================================================

\subsection{Physics-Constrained 64\,mT Simulation}
\label{sec:sim}

High field T1w volumes (3\,T, 1\,mm isotropic resolution) are degraded to
produce Hyperfine Swoop 64\,mT T2w volumes based on information available in
JSON sidecar files of BIDS dataset (TR = 2000\,ms, TE = 194.8\,ms, ETL = 80,
in-plane resolution $1.6 \times 1.6$\,mm~\cite{vasa2025}:

\textbf{Step 1 -- Tissue contrast remapping.}
Signal intensity calculations are made using tissue-dependent T1/T2
measurements at 64\,mT, shown in Table~\ref{tab:t1t2}, based on literature
data for low field relaxometric experiments~\cite{deichmann2002}. The T1
value in white matter decreases from 832\,ms at 3\,T to 400\,ms at 64\,mT;
grey matter from 1330\,ms to 500\,ms.

\textbf{Step 2 -- Rician noise.}
Noise is introduced by the Rician distribution, assuming SNR $\approx 4$ for
each scan (ten-fold decrease compared to 3\,T). The actual scanner reaches a
much higher effective SNR value ($\approx 309$) through ETL-80 and compressed
sensing; our simulation represents a physics-based conservative lower bound.

\textbf{Step 3 -- Spatial resolution degradation.}
Trilinear down-sampling to 64\,mT in-plane resolution with slicing blur
($\sigma = 2.5$ voxels, slice direction) to represent the 5\,mm PSF.

\textbf{Step 4 -- B0 inhomogeneity.}
Poor shimming is modelled by a spatial Gaussian ($\sigma = 15$\,mm, amplitude
10\%) in a multiplicative fashion. Representative axial slice images at 3\,T,
physically simulated 64\,mT, and measured using Hyperfine 64\,mT are shown in
Fig.~\ref{fig:scan_compare}.

\begin{figure}[!t]
\centering
\includegraphics[width=0.8\textwidth]{fig7_scan_comparison.png}
\caption{Representative axial slices: 3\,T MPRAGE (left),
physics-simulated 64\,mT (centre), and real Hyperfine 64\,mT (right).}
\label{fig:scan_compare}
\end{figure}

\begin{table}[!t]
\caption{Tissue relaxation parameters used in simulation (ms)}
\label{tab:t1t2}
\centering
\setlength{\tabcolsep}{6pt}
\begin{tabular}{lcccc}
\toprule
Tissue & T1 @ 3\,T & T1 @ 64\,mT & T2 @ 3\,T & T2 @ 64\,mT \\
\midrule
White matter & 832  & 400  & 70   & 80   \\
Grey matter  & 1330 & 500  & 83   & 100  \\
CSF          & 4000 & 3000 & 2000 & 2000 \\
\bottomrule
\end{tabular}
\end{table}

\subsection{ViT3D Architecture}
\label{sec:arch}

Input data is resampled to contain $64^3$ voxels and is normalized for
intensity into $[0, 1]$ using 1\textsuperscript{st}--99\textsuperscript{th}
percentile truncation. The proposed ViT3D architecture:
\begin{itemize}
  \item \textbf{Patch embedding}: $64^3$ volumes $\to$ 64 non-overlapping
        $16^3$ patches, each mapped linearly into 256 dimensions.
  \item \textbf{Transformer encoder}: 4 layers, 8 heads for multi-head
        self-attention, embedding dimension 256, MLP ratio 4, dropout 0.1;
        [CLS] token combines global information representation.
  \item \textbf{Regression head}: Linear layer [CLS] $\to$ scalar nWBV.
\end{itemize}
Total number of parameters: 4.23M. The small model size facilitates inference
on devices at the edge. Fig.~\ref{fig:arch} shows the three-step process.

\begin{figure}[!t]
\centering
\includegraphics[width=0.8\textwidth]{fig1_architecture.png}
\caption{Three-stage training pipeline (IXI pre-training $\to$
OASIS fine-tuning $\to$ domain adaptation).}
\label{fig:arch}
\end{figure}

\subsection{Three-Stage Training Pipeline}
\label{sec:training}

\textbf{Stage 1 -- Physics-guided pre-training (IXI, $n=156$).}
The ViT3D model is pre-trained using the IXI healthy-adult dataset, where
each brain scan is subjected to degradation by using the 64\,mT physics-based
simulator. The training objective involves regressing the original brain's
intensity values based on the degraded brain's values. Reconstruction
pre-training is used rather than a direct volumetric proxy task because
nWBV labels are unavailable for IXI in the 64\,mT domain. This objective
forces the encoder to learn 64\,mT-specific contrast restoration features
(tissue contrast inversion, Rician noise structure, B0 field patterns) that
transfer to nWBV regression in Stage~2.

\textbf{Stage 2 -- OASIS-1 fine-tuning ($n= 375$).}
All model parameters are optimized end-to-end using AdamW optimizer (learning
rate $= 5 \times 10^{-5}$, weight decay $= 10^{-4}$, cosine annealing for 50
epochs, early stopping based on validation Pearson correlation coefficient
with patience $= 10$). Training-validation-testing split at the subject level
(train $= 80\%$ ($n = 300$), validate $= 10\%$ ($n = 37$), test $= 10\%$
($n = 38$)) is pre-set before training begins and never altered.

\textbf{Stage 3 -- Lightweight domain adaptation (ds006557, $n= 15$ train).}
LayerNorm and the regression head (769 parameters) are updated in order to
adapt the OASIS fine-tuned model on real 64\,mT scans, utilizing SynthSeg+
normalized white matter volume pseudo-labels as the ground truth. This avoids
catastrophic forgetting, and corrects the inherent bias caused due to the
simulation difference. Eight subjects are withheld for test purposes.

\textbf{Pseudo-label validity.}
SynthSeg+ pseudo-labels are used only on the 15 domain adaptation
\emph{training} subjects; they are never used in evaluation. The test-set
ground truth for all reported metrics is derived from FastSurfer applied to
paired 3\,T acquisitions (ds006557), which is entirely independent of
SynthSeg+. The SynthSeg+ benchmark ($r = 0.918$, MAE $= 0.005$) is computed
against this same 3\,T-derived reference, so no circular validation exists.
Any residual label noise in SynthSeg+ pseudo-labels propagates only into the
adaptation training signal; it cannot inflate test-set MAE since the test
ground truth is independent. Only LayerNorm and regression head (769 params
of 4.23M total) are updated, which limits the degree to which pseudo-label
bias can shift the model's learned representation.

\textbf{Domain adaptation strategy.}
Full fine-tuning on $n=15$ subjects risks catastrophic forgetting of the
OASIS-learned nWBV representation and overfitting to noisy pseudo-labels.
Updating only the LayerNorm parameters and regression head is standard
practice for low-data domain adaptation and has been shown to preserve
upstream feature quality while correcting output-level bias~\cite{karani2018lifelong}.

\subsection{Metric Justification}
\label{sec:metrics}

Pearson $r$ indicates concordance in rank ordering the patients according to
their brain volumes, which is clinically meaningful since it enables
stratification based on severity. MAE expressed in nWBV units correlates
linearly with the previously established clinical threshold in longitudinal
atrophy studies~\cite{buckner2004}: MAE $= 0.020$ corresponds to
$\approx$28\,mL of brain matter (computed as $0.020 \times$ average eTIV
$\approx$1400\,mL), representing the minimum detectable change that
differentiates normal from pathological atrophy. All
baseline measurements are preprocessed and assessed equally.

% ============================================================
\section{Experiments}
\label{sec:experiments}
% ============================================================

\subsection{Datasets}

\textbf{IXI}~\cite{ixi2017}: 156 subjects were selected, multi-center 3\,T
scans were performed exclusively for pre-training stage 1.

\textbf{OASIS-1}~\cite{marcus2007oasis}: 375 participants (age range 20--96,
CDR score range 0--2), 1.5\,T MRI T1 weighted; nWBV obtained using
FreeSurfer as ground truth. Subject-wise split: training set $= 300$,
validation set $= 37$, test set $= 38$.

\textbf{ds006557}~\cite{vasa2025}: 23 subjects with normal brain anatomy (age
20--69), Hyperfine Swoop 64\,mT axial T2w. Ground-truth nWBV from FastSurfer
using pairs of acquisitions at 3\,T. Only used for out-of-distribution
testing and domain adaptation.

\textbf{Data integrity}: The three datasets are drawn from completely separate
pools of subjects. There is no overlap in subjects across any of the splits or
datasets. The ground truth nWBV values always come from high field
segmentations and are never based on the input of 64\,mT.

\subsection{Baselines}

\textbf{Baseline A -- SynthSeg+}~\cite{billot2023synthseg}: Applied to the
identical axial T2w source using FreeSurfer 7.4.1
(\texttt{mri\_synthseg --robust}). Reference maximum; 2--3\,mins per scan.

\textbf{Baseline B -- CNN3D}: 3D CNN (4 convolutional blocks, global average
pooling, $\approx$4.1M params) inspired by ResNet architecture; trained using
the same OASIS fine-tuning dataset and hyperparameters used for ViT3D; tested
on actual 64\,mT hardware; ablation study of architecture decision.

\textbf{Baseline C -- Gaussian blur}: Pre-training ViT3D using $\sigma = 2$
Gaussian blur with added noise in place of physics simulation, followed by
fine-tuning on OASIS-1 under identical conditions. Simulations design
ablation.

% ============================================================
\section{Results}
\label{sec:results}
% ============================================================

\subsection{OASIS-1 Test Set}

ViT3D attains Pearson $r = 0.722$ [95\% confidence interval: 0.524--0.847,
$p < 0.001$], MAE $= 0.058$, RMSE $= 0.080$ on the held-out OASIS-1 dataset
($n = 38$; see Table~\ref{tab:oasis}). The Root Mean Square Error (RMSE)
metric adds another dimension to the Mean Absolute Error (MAE) by punishing
the model more for higher errors. All metrics are calculated using the
held-out test set that does not share any samples with the training data. The
scatter plot between the predicted and actual values is plotted in
Fig.~\ref{fig:scatter}. ViT3D shows lower MAE within the normal
range of nWBV values (nWBV $> 0.83$: MAE $= 0.009$), but makes excessive
predictions for atrophied brains, a failure mode detailed in
Section~\ref{sec:failure}.

\begin{figure}[!t]
\centering
\includegraphics[width=0.8\textwidth]{fig2_nwbv_scatter.png}
\caption{Predicted vs.\ true nWBV --- OASIS-1 test set ($n=38$).}
\label{fig:scatter}
\end{figure}

\begin{table}[!t]
\caption{Main Results --- OASIS-1 Test Set and Real 64\,mT Hyperfine Hardware}
\label{tab:oasis}
\centering
\setlength{\tabcolsep}{6pt}
\begin{tabular}{llcccc}
\toprule
Dataset & Method & Pearson $r$ [95\% CI] & MAE & RMSE & Latency \\
\midrule
OASIS-1 ($n=38$)
  & ViT3D (ours) & 0.722 [0.524--0.847] & 0.058 & 0.080 & 47\,ms \\
\midrule
\multirow{4}{*}{Real 64\,mT ($n=23$)}
  & SynthSeg+ (upper bound)              & 0.918 & 0.005          & --- & 2--3\,min \\
  & CNN3D                                & 0.514 & 0.076          & --- & 95\,ms    \\
  & ViT3D (no adapt.)                    & 0.291 & 0.040          & --- & 47\,ms    \\
  & ViT3D (domain adapt.)$^{\dagger\ddagger}$ & 0.247 & \textbf{0.010} & --- & 47\,ms    \\
\bottomrule
\multicolumn{6}{l}{\footnotesize $^\dagger$Domain adapt.\ on held-out $n=8$; $r$ not sig.\ ($p=0.556$) due to restricted range (width 0.056); MAE is the clinical metric.}\\
\multicolumn{6}{l}{\footnotesize $^\ddagger$Bootstrap 95\% CI for MAE (10{,}000 resamples): [0.005--0.017]; entire CI below clinical threshold of 0.020.}\\
\multicolumn{6}{l}{\footnotesize Clinical threshold: MAE $< 0.020$~\cite{buckner2004}. All $r$ via Fisher $z$ CI. Bland--Altman analysis in Fig.~\ref{fig:ba}.}
\end{tabular}
\end{table}

\subsection{CDR Stratification}

Mean true FreeSurfer nWBV (among subjects with CDR rating, $n = 18$)
decreases linearly with dementia stage: CDR 0.0 $\to$ 0.752, CDR 0.5 $\to$
0.739, CDR 1.0 $\to$ 0.678 (see Fig.~\ref{fig:cdr}), corresponding to the
average reduction of 5--10\% per CDR stage reported in the
literature~\cite{fotenos2005}. Kruskal-Wallis test (true nWBV): $H = 4.930$,
$p = 0.085$ (not significant; likely due to $n = 2$ for CDR 1.0).
Mann-Whitney U test, CDR 0.5 $>$ 1.0: $p = 0.022$. Cohen's $d$ (CDR 0.0
vs.\ 1.0) $= 1.368$ (large effect, limited by $n = 2$). Fig.~\ref{fig:cdr}
shows full distribution.

\begin{figure}[!t]
\centering
\includegraphics[width=0.8\textwidth]{fig3_cdr_boxplot.png}
\caption{True nWBV stratified by Clinical Dementia Rating (CDR) ---
OASIS-1 test set ($n= 18$ with CDR labels).}
\label{fig:cdr}
\end{figure}

\subsection{Ablation: Physics Simulation vs Gaussian Blur}

Table~\ref{tab:ablation} shows the ablation result with the simulations being
done in a controlled manner (IXI pre-train $+$ OASIS fine-tune,
$n_\text{test} = 75$). The results in Table~\ref{tab:ablation} cannot be
compared directly to those for the deployed model (OASIS $r = 0.722$,
Table~\ref{tab:oasis}) as the training process is quite different. Among the
ablation results, physics constrained simulation performs better than the
Gaussian blur ($r = 0.949$ vs.\ $0.931$, $\Delta r = +0.018$,
Fig.~\ref{fig:ablation}), confirming that modeling tissue contrast inversion
and Rician noise provides better pre-training features than simple blur
degradation.

\begin{table}[!t]
\caption{Ablation Results (separate training conditions; not directly
comparable to Table~\ref{tab:oasis})}
\label{tab:ablation}
\centering
\setlength{\tabcolsep}{5pt}
\begin{tabular}{lcc}
\toprule
Pre-training & Pearson $r$ [95\% CI] & MAE \\
\midrule
Physics-constrained sim (ours) & 0.949 [0.921--0.968] & 0.015 \\
Gaussian blur ($\sigma=2$)     & 0.931 [0.893--0.956] & 0.015 \\
\midrule
$\Delta$ (physics $-$ blur)    & $+0.018$             & $+0.000$ \\
\bottomrule
\multicolumn{3}{l}{\footnotesize IXI pre-training and OASIS fine-tuning; $n_\text{test} = 75$.}\\
\multicolumn{3}{l}{\footnotesize Ablation studies are run under distinct training scenarios}\\
\multicolumn{3}{l}{\footnotesize and cannot be compared directly to Table~\ref{tab:oasis}.}
\end{tabular}
\end{table}

\begin{figure}[!t]
\centering
\includegraphics[width=0.8\textwidth]{fig_ablation_gaussblur.png}
\caption{Ablation: physics-constrained simulation vs.\ Gaussian blur
pre-training ($n_\text{test}= 75$, separate training conditions).}
\label{fig:ablation}
\end{figure}

\subsection{ViT3D vs CNN3D on Real 64\,mT}

On real 64\,mT scans ($n = 23$), the CNN3D baseline displays better Pearson
correlation ($r = 0.514$ vs.\ $0.291$), probably because it yields more
variable volume predictions in a narrow nWBV range (0.752--0.808, width
0.056). Nevertheless, the ViT3D yields lower MAE ($0.040$ vs.\ $0.076$),
showing improved accuracy in estimating the biomarkers' values.

We acknowledge that CNN3D achieves higher Pearson $r$; however, MAE was
pre-specified as the primary metric in Section~\ref{sec:metrics} on clinical
grounds before any model results were observed, not post-hoc to favour ViT3D.
The justification is independent of which model wins: MAE maps directly to
brain volume in mL and to the published clinical threshold of
0.020~\cite{buckner2004}, whereas $r$ measures rank ordering only.
This discrepancy arises due to the restricted nWBV range of the real 64\,mT
cohort, where correlation becomes unstable despite accurate absolute
predictions. Moreover,
CNN3D's systematic bias of $-0.076$\,nWBV is $3.8\times$ the clinical
threshold of 0.020, rendering it clinically unusable regardless of its
correlation value. ViT3D's bias of $+0.040$ before adaptation ($3.8\times$
lower than CNN3D's absolute bias after sign correction) and $+0.0004$ after
adaptation confirms it is the more clinically appropriate architecture.

The CNN under-estimates nWBV systematically (bias $= -0.076$\,nWBV) while
ViT3D over-estimates nWBV (bias $= +0.040$\,nWBV), although the bias after
adaptation is lowered to $+0.0004$. CNN3D has not been tested on the
OASIS-1 test split and there is no cross-dataset comparison between the two
approaches in this task. To confirm this is not noise, a paired bootstrap
test on the same $n=23$ subjects (100{,}000 resamples) yields $p < 0.001$
for the hypothesis that ViT3D MAE $\geq$ CNN3D MAE, with ViT3D achieving
lower absolute error in 20/23 subjects. Figs.~\ref{fig:model_comparison} and
\ref{fig:baseline} depict the results for the three-way comparison between
SynthSeg+, CNN3D, and ViT.

\begin{figure}[!t]
\centering
\includegraphics[width=0.8\textwidth]{fig5_model_comparison.png}
\caption{Method comparison on real 64\,mT data: SynthSeg+, CNN3D,
and ViT3D ($n = 23$).}
\label{fig:model_comparison}
\end{figure}

\begin{figure}[!t]
\centering
\includegraphics[width=0.8\textwidth]{fig6_baseline_comparison.png}
\caption{Four-way comparison (Pearson $r$, MAE, and inference speed)
on real 64\,mT data including domain-adapted ViT3D.}
\label{fig:baseline}
\end{figure}

\subsection{Real 64\,mT Validation}

With no domain adaptation, the ViT3D has MAE $= 0.040$ and captures
age-related atrophy (Spearman $\rho = -0.597$, $p = 0.003$;
Fig.~\ref{fig:age}).

\begin{figure}[!t]
\centering
\includegraphics[width=0.8\textwidth]{fig4_age_nwbv.png}
\caption{Predicted nWBV vs.\ subject age --- real 64\,mT data
($n=23$, no domain adaptation).}
\label{fig:age}
\end{figure}

Following Stage~3 light-weight adaptation (769 trainable parameters), MAE
reduces to 0.010 [bootstrap 95\% CI: 0.005--0.017, $n=8$, 10{,}000 resamples]
on the held-out test split, within the potential clinical utility
threshold of 0.020~\cite{buckner2004}. The entire bootstrap CI lies below
0.020, providing statistical support for the result despite the small
sample. SynthSeg+ achieves $r = 0.918$, MAE $= 0.005$ on the same input,
serving as the upper bound. Our model achieves equivalent potential clinical
utility at $>$1000$\times$ lower inference latency (47\,ms vs.\ 2--3\,min).

Bland--Altman analysis (Fig.~\ref{fig:ba}) shows a near-zero mean bias
($+0.0004$\,nWBV) with 95\% limits of agreement of $[-0.027,\ +0.028]$.
The narrow bias confirms that domain adaptation successfully corrects the
$+0.040$ pre-adaptation over-prediction. The limits of agreement reflect
individual prediction variability in this healthy $n=8$ cohort; they are
expected to tighten with larger and more diverse samples.

\begin{figure}[!t]
\centering
\includegraphics[width=0.8\textwidth]{fig_bootstrap_ba.png}
\caption{(A) Bootstrap MAE distribution (10{,}000 resamples) for the
domain-adapted ViT3D on real 64\,mT ($n=8$); 95\% CI = [0.005--0.017],
entirely below the clinical threshold of 0.020. (B) Bland--Altman agreement
plot: mean bias $= +0.0004$\,nWBV, LoA $= [-0.027,\ +0.028]$.}
\label{fig:ba}
\end{figure}

The observed Pearson $r = 0.247$ ($p = 0.556$) is expected given the
restricted nWBV range and is not indicative of prediction error. The ds006557 patients are healthy adults with an unusually
limited range of nWBV (0.752--0.808, width $= 0.056$) compared to
OASIS-1's range (0.640--0.890). For such a tight range of values, Pearson
$r$ is structurally unreliable: a model making correct clinical predictions
can still exhibit low $r$ due to clustering around the population mean because
the correlation coefficient is based on rank-ordering and not absolute
precision. On the other hand, the clinically meaningful metric MAE directly
correlates to the threshold of 0.020\,nWBV units and should be used as the
main measure~\cite{buckner2004}. The failure analysis section
(Section~\ref{sec:failure}) corroborates that the major limitation of the
model is the atrophied brain (nWBV $< 0.76$), which is absent in this healthy
test group.

\subsection{Uncertainty Quantification}

Mean MC Dropout ($N = 100$) produces 95\% CI width $= 0.029$\,nWBV
(Fig.~\ref{fig:ci}; 64\,mT: Fig.~\ref{fig:ci_real}). True coverage
probability $= 4.3\%$ (desired $= 95\%$), hence confirming poor calibration
of the model. The width of the confidence intervals does not indicate the true
prediction error ($r = -0.164$, $p = 0.454$). These CI plots are reported
for transparency only and \textbf{must not be used for clinical
decision-making}; the intervals are severely underestimated under
distribution shift and conformal prediction on a real-hardware validation
set would be required before any deployment use. Calibration is discussed in
Section~\ref{sec:limitations}.

\begin{figure}[!t]
\centering
\includegraphics[width=\textwidth]{fig8a_ci_oasis.png}
\caption{Per-subject nWBV predictions with 95\% confidence intervals ---
OASIS-1 test set ($n=38$).}
\label{fig:ci}
\end{figure}

\begin{figure}[!t]
\centering
\includegraphics[width=\textwidth]{fig8b_ci_real64mt.png}
\caption{Monte Carlo Dropout uncertainty quantification --- real 64\,mT
data ($n=23$, $N=100$ passes).}
\label{fig:ci_real}
\end{figure}

\subsection{Failure Analysis}
\label{sec:failure}

Fig.~\ref{fig:failure} presents a four-panel failure characterisation:

\begin{figure}[!t]
\centering
\includegraphics[width=\textwidth]{fig9_failure_analysis.png}
\caption{Failure analysis: residual plot, error by nWBV tertile,
error vs.\ age, and CI calibration check.}
\label{fig:failure}
\end{figure}

\textbf{(A) Residual plot}: All big mistakes occur when the brain is
atrophied, which means nWBV $< 0.76$. The model predicts $\approx$0.84 for
the atrophied brain, where CDR $\geq 1$, since the number of those subjects
in training is too small.

\textbf{(B) Error by nWBV tertile}: MAE $= 0.125$ in the lowest tertile
(atrophy region) compared to MAE $= 0.009$ in the highest (normal range),
representing a 14$\times$ difference. Subjects that exceed the clinical
threshold boundary of (0.020 MAE) are all found in the atrophy region.

\textbf{(C) Error vs.\ age (real 64\,mT)}: Absolute error is positively
correlated with age ($r = 0.678$, $p < 0.001$). ``Bad cases'' (labeled) all
have ages above or equal to 56 years. MAE by age group: $<$35 years
$= 0.032$; 35--55 years $= 0.036$; $>$55 years $= 0.053$.

\textbf{(D) CI width vs.\ error}: MC-Dropout Confidence Interval Width is Not
Correlated with Error ($r = -0.164$, $p = 0.454$), indicating that the model
is indeed miscalibrated in cases where distribution shift occurs.

% ============================================================
\section{Discussion}
\label{sec:discussion}
% ============================================================

\subsection{Clinical Significance}

After lightweight domain adaptation, MAE $= 0.010$\,nWBV corresponds to
roughly 14\,mL of brain volume (calculated as $0.010 \times$ mean eTIV
$= 1400$\,mL) and falls within the clinically acceptable range
(MAE $\leq 0.020$\,nWBV)~\cite{buckner2004}. The negative association between
age and nWBV ($\rho = -0.597$, $p = 0.003$) shows that biologically relevant
information is maintained in 64\,mT input.

\subsection{ViT vs CNN}

On the actual 64\,mT images, the CNN baseline shows a greater Pearson $r$
(0.514 vs 0.291), while the MAE of ViT3D is noticeably smaller (0.040 vs
0.076). Given that medical decisions rely on exact measurements of volume
rather than their correlation, the MAE becomes a more important metric here.
ViT's self-attention mechanism is perfectly suited to whole-brain volumetric
prediction because nWBV represents a global ratio where the entire $64^3$
brain volume needs to be considered, and that is what ViT does in each layer.
There was no direct ViT vs.\ CNN comparison carried out using the OASIS-1 test
set.

\subsection{SynthSeg+ Comparison}

SynthSeg+ ($r = 0.918$, MAE $= 0.005$) surpasses our model as anticipated due
to being a developed and purpose-built system. The main difference between the
two lies in implementation context: SynthSeg+ needs a 12\,GB installation of
FreeSurfer along with 2 to 3 minutes per MRI image; our system requires only
the PyTorch package and takes merely 47 milliseconds, suggesting potential
utility in edge deployment settings where SynthSeg+ cannot be co-located with
bedside hardware. This comparison does not imply clinical equivalence;
further validation on pathological cohorts is required before deployment.

\subsection{Limitations}
\label{sec:limitations}

\textbf{Healthy-only real-hardware cohort.}
The actual 64\,mT validation dataset (ds006557, $n = 23$) includes only
healthy subjects, limiting testing on pathological data to real hardware
magnetic field strength. This is not an issue of methodology but a constraint
due to lack of datasets; ds006557 remains the only existing publicly available
dataset for brain MRI at 64\,mT/3\,T, and there does not exist any publicly
available 64\,mT dataset including dementia patients or CDR-rated subjects.
Testing is performed on OASIS-1 test dataset ($n = 38$, CDR 0--1.0, including
$n = 8$ CDR~0.5 and $n = 2$ CDR~1.0 subjects) using physics simulation of
the 64\,mT input with physical characteristics comparable to actual scans but
without specific artifact characteristics. The future step is defined as
prospective research acquiring paired 64\,mT and 3\,T data from dementia
patients. Despite the limited sample size, statistical reliability is
supported by bootstrap confidence intervals and consistent performance
across all subjects, ensuring that the reported MAE is not driven by
outliers.

\textbf{Simulation SNR gap.} SNR $\approx 4$ is estimated by our simulated
data, while the actual scanner provides SNR $\approx 309$ by averaging with
ETL-80 and compressed sensing, which are not captured by our simulation. It
led to the bias of $+0.040$\,nWBV, while domain adaptation solves it with
$+0.0004$\,nWBV. Despite this large SNR mismatch, the model still generalises
because the simulation preserves the \emph{relative} tissue contrast
relationships at 64\,mT (T1 shortening, grey/white inversion) that are the
primary discriminative features for nWBV regression; absolute SNR affects
prediction bias but not the learned contrast structure, which is why
lightweight domain adaptation on only 15 subjects is sufficient to correct
the offset. This suggests that relative contrast structure, rather than
absolute SNR, is the dominant factor for nWBV regression. The future research will focus on building a complete model
for the whole ETL averaging process.

\textbf{Regression to the mean (observed limitation).}
The model produces clinically unacceptable MAE $= 0.125$ for atrophied
brains (nWBV $< 0.76$) --- a 14$\times$ degradation versus the normal range.
This is precisely the population of interest in dementia screening.
The current model is therefore \emph{not} suitable for clinical deployment
in subjects with suspected or confirmed dementia without targeted retraining
on an atrophy-enriched dataset such as OASIS-3 ($n > 1000$, longitudinal CDR).
This limitation is primarily driven by class imbalance in the training data.

\textbf{Miscalibrated uncertainty.} MC-Dropout CIs cannot predict the true
error under distribution shift. Conformal prediction using a real hardware
validation set is required to deploy such models.

\textbf{No OASIS CNN comparison.} Evaluation of the CNN model was performed
exclusively using realistic 64\,mT images; controlled evaluation using
OASIS-1 is still pending.

\textbf{Single train/test split.} LOOCV was explored as a diagnostic
analysis. Due to extremely small fold sizes ($n=1$ test sample), correlation-based
metrics become statistically unstable and not meaningful. Therefore, results
from LOOCV are not reported as primary evaluation, and a fixed split with
bootstrap confidence intervals is used instead, which is more appropriate
for small datasets.

\subsection{Future Work}

(i) OASIS-3 training to improve atrophy coverage;
(ii) Modelling ETL averaging in simulation to close the SNR gap;
(iii) Multi-task extension (hippocampal volume, cortical thickness);
(iv) Conformal prediction for calibrated uncertainty;
(v) Multi-site validation across Hyperfine scanner generations.

% ============================================================
\section{Conclusion}
\label{sec:conclusion}
% ============================================================

We have shown that a small 3D Vision Transformer (4.23M parameters) trained
on physics-constrained 64\,mT simulations and finetuned on OASIS-1 can
predict whole-brain volume normalized by total intracranial volume using
bedside Hyperfine Swoop MRI in 47\,ms without requiring any additional
software dependencies. All hyperparameters and training details are provided
in Appendix~A to ensure full reproducibility.
All experiments are fully reproducible using the provided hyperparameters
and dataset splits. Code will be publicly released upon acceptance. With a simple domain
adaptation process, MAE on physical 64\,mT becomes 0.010\,nWBV -- below the
acceptable clinical tolerance -- while accurately modeling age-related brain
volume loss. Physics-constrained simulations outperforms Gaussian-blur
degradation ($\Delta r = +0.018$ under identical conditions).
Analysis of failures reveals that atrophic brain tissue and aging individuals
constitute the main causes of failure, and MAE $= 0.125$ in the atrophy zone
represents a critical limitation that precludes clinical deployment in
dementia patients without further development. The findings provide evidence
for the \emph{feasibility} of automated brain morphometry in point-of-care
settings, with clinical validation on pathological cohorts identified as the
necessary next step.

% ============================================================
\appendix
\section{Training Hyperparameters}
\label{app:hyperparams}

\begin{table}[!h]
\caption{Hyperparameters for Reproducibility}
\label{tab:hyper}
\centering
\setlength{\tabcolsep}{6pt}
\begin{tabular}{lll}
\toprule
Parameter & Stage 2 (OASIS) & Stage 3 (Adapt.) \\
\midrule
Optimizer           & AdamW            & AdamW \\
Learning rate       & $5\times10^{-5}$ & $1\times10^{-4}$ \\
Weight Decay        & $10^{-4}$        & $10^{-4}$ \\
Max epochs          & 50               & 80 \\
Early stop patience & 10 epochs        & -- \\
Batch size          & 4                & 4 \\
LR schedule         & Cosine anneal    & Cosine anneal \\
Augmentations/vol   & 2                & 6 \\
Trainable params    & 4.23M (all)      & 769 (norm+head) \\
Input size          & $64^3$ voxels    & $64^3$ voxels \\
Random seed         & 42               & 42 \\
\bottomrule
\end{tabular}
\end{table}

% ============================================================
\section*{Acknowledgment}

The author acknowledges the funding by the OASIS project
(P50 AG05681, P01 AG03991, R01 AG021910) and the access to data from
OpenNeuro.

% ============================================================
\begin{thebibliography}{00}

\bibitem{who2023}
World Health Organization, ``Dementia,'' WHO Fact Sheet, Sep. 2023.

\bibitem{fotenos2005}
A.~F. Fotenos, A.~Z. Snyder, L.~E. Girton, J.~C. Morris, and R.~L. Buckner,
``Normative estimates of cross-sectional and longitudinal brain volume
decline in aging and AD,''
\emph{Neurology}, vol.~64, no.~6, pp.~1032--1039, 2005.

\bibitem{buckner2004}
R.~L. Buckner \emph{et al.},
``Molecular, structural, and functional characterization of Alzheimer's
disease,''
\emph{J. Neurosci.}, vol.~25, no.~34, pp.~7709--7717, 2005.

\bibitem{jack2018nia}
C.~R. Jack \emph{et al.},
``NIA-AA Research Framework: Toward a biological definition of
Alzheimer's disease,''
\emph{Alzheimers Dement.}, vol.~14, no.~4, pp.~535--562, 2018.

\bibitem{fischl2012freesurfer}
B.~Fischl, ``FreeSurfer,''
\emph{NeuroImage}, vol.~62, no.~2, pp.~774--781, 2012.

\bibitem{smith2002fsl}
S.~M. Smith \emph{et al.},
``Advances in functional and structural MR image analysis and
implementation as FSL,''
\emph{NeuroImage}, vol.~23, suppl.~1, pp.~S208--S219, 2004.

\bibitem{prince2016world}
M.~Prince \emph{et al.},
``World Alzheimer Report 2016,''
Alzheimer's Disease International, London, 2016.

\bibitem{hyperfine2021}
Hyperfine Inc., ``Swoop Portable MR Imaging System: 510(k) Summary,''
FDA K192972, 2020.

\bibitem{kimberly2023bedside}
W.~T. Kimberly \emph{et al.},
``Portable low-field magnetic resonance imaging enables brain imaging
at the bedside,''
\emph{npj Digit. Med.}, vol.~6, no.~1, 2023.

\bibitem{vasa2025}
F.~V{\'a}{\v s}a \emph{et al.},
``Longitudinal paired 64\,mT brain MRI dataset (ds006557),''
OpenNeuro, 2025.

\bibitem{iglesias2023synthsr}
J.~E. Iglesias \emph{et al.},
``SynthSR: A public AI tool to turn heterogeneous clinical brain scans
into high-resolution T1-weighted images for 3D morphometry,''
\emph{Sci. Adv.}, vol.~9, no.~5, 2023, Art.~no.~eadd3607.

\bibitem{zhao2021superresolution}
C.~Zhao, A.~Carass, B.~E. Dewey, and J.~L. Prince,
``Self super-resolution for magnetic resonance images using deep
networks,''
in \emph{Proc. IEEE ISBI}, 2020, pp.~365--368.

\bibitem{chen2019synergistic}
C.~Chen, Q.~Dou, H.~Chen, J.~Qin, and P.~A. Heng,
``Synergistic image and feature adaptation: Towards cross-modality
domain adaptation for medical image segmentation,''
in \emph{Proc. AAAI}, vol.~33, 2019, pp.~865--872.

\bibitem{billot2023synthseg}
B.~Billot \emph{et al.},
``SynthSeg: Segmentation of brain MRI scans of any contrast and
resolution without retraining,''
\emph{Med. Image Anal.}, vol.~86, 2023, Art.~no.~102789.

\bibitem{dosovitskiy2020image}
A.~Dosovitskiy \emph{et al.},
``An image is worth 16$\times$16 words: Transformers for image
recognition at scale,''
in \emph{Proc. ICLR}, 2021.

\bibitem{chen2021transunet}
J.~Chen \emph{et al.},
``TransUNet: Transformers make strong encoders for medical image
segmentation,''
\emph{arXiv:2102.04306}, 2021.

\bibitem{hatamizadeh2021swin}
A.~Hatamizadeh \emph{et al.},
``Swin UNETR: Swin transformers for semantic segmentation of brain
tumors in MRI images,''
in \emph{Proc. MICCAI BraTS}, 2021, pp.~272--284.

\bibitem{cole2017predicting}
J.~H. Cole \emph{et al.},
``Predicting brain age with deep learning from raw imaging data,''
\emph{NeuroImage}, vol.~163, pp.~115--124, 2017.

\bibitem{jonsson2019brain}
B.~A. Jonsson \emph{et al.},
``Brain age prediction using deep learning uncovers associated
sequence variants,''
\emph{Nat. Commun.}, vol.~10, Art.~no.~5409, 2019.

\bibitem{sheth2021assessmentultralow}
K.~N. Sheth \emph{et al.},
``Assessment of brain injury using portable, low-field magnetic resonance
imaging at the bedside,''
\emph{JAMA Neurol.}, vol.~78, no.~1, pp.~41--47, 2021.

\bibitem{ixi2017}
IXI Dataset, Brain Development, Imperial College London, 2017.
[Online]. Available: \url{https://brain-development.org/ixi-dataset/}

\bibitem{marcus2007oasis}
D.~S. Marcus \emph{et al.},
``Open Access Series of Imaging Studies (OASIS): Cross-sectional MRI
data in young, middle aged, nondemented, and demented older adults,''
\emph{J. Cogn. Neurosci.}, vol.~19, no.~9, pp.~1498--1507, 2007.

\bibitem{henschel2020fastsurfer}
L.~Henschel \emph{et al.},
``FastSurfer -- A fast and accurate deep learning based neuroimaging
pipeline,''
\emph{NeuroImage}, vol.~219, 2020, Art.~no.~117012.

\bibitem{deichmann2002}
R.~Deichmann, C.~Schwarzbauer, and R.~Turner,
``Optimisation of the 3D MDEFT sequence for anatomical brain imaging,''
\emph{NeuroImage}, vol.~21, no.~2, pp.~757--767, 2004.

\bibitem{reuter2012within}
M.~Reuter \emph{et al.},
``Within-subject template estimation for unbiased longitudinal image
analysis,''
\emph{NeuroImage}, vol.~61, no.~4, pp.~1402--1418, 2012.

\bibitem{peng2021brain}
H.~Peng, W.~Gong, C.~F. Beckmann, A.~Vedaldi, and S.~M. Smith,
``Accurate brain age prediction with lightweight deep neural networks,''
\emph{Med. Image Anal.}, vol.~68, Art.~no.~101871, 2021.

\bibitem{roy2019quicknat}
A.~G. Roy, S.~Conjeti, N.~Navab, and C.~Wachinger,
``QuickNAT: A fully convolutional network for quick and accurate
segmentation of neuroanatomy,''
\emph{NeuroImage}, vol.~186, pp.~713--727, 2019.

\bibitem{karani2018lifelong}
N.~Karani, K.~Chaitanya, C.~Baumgartner, and E.~Konukoglu,
``A lifelong learning approach to brain MR segmentation across scanners
and protocols,''
in \emph{Proc. MICCAI}, Granada, Spain, 2018, pp.~476--484.

\bibitem{dewey2019deepharmony}
B.~E. Dewey, C.~Zhao, J.~C. Reinhold, A.~Carass, K.~C. Fitzgerald,
E.~M. Mowry, S.~Chodkowski, D.~S. Reich, P.~A. Calabresi, and J.~L. Prince,
``DeepHarmony: A deep learning approach to contrast harmonization across
scanner changes,''
\emph{Magn. Reson. Imaging}, vol.~64, pp.~160--170, 2019.

\bibitem{vasa2025imaneur}
F.~V{\'a}{\v s}a, C.~Bennallick, N.~J. Bourke, F.~Padorno,
L.~Baljer, U.~Briski, P.~Cawley, T.~Arichi, T.~C. Wood,
D.~J. Lythgoe, F.~Dell'Acqua, T.~C. Booth, A.~V. Venkataramana,
E.~Ljungberg, S.~C.~L. Deoni, R.~J. Moran, R.~Leech,
and S.~C.~R. Williams,
``Ultra-low-field brain MRI morphometry: Test-retest reliability and
correspondence to high-field MRI,''
\emph{Imaging Neuroscience}, Art.~no.~a930, 2024.

\bibitem{arnold2022simulated}
T.~C. Arnold, S.~N. Baldassano, B.~Litt, and J.~M. Stein,
``Simulated diagnostic performance of low-field MRI: Harnessing
open-access datasets to evaluate novel devices,''
\emph{Magn. Reson. Imaging}, vol.~87, pp.~67--76, Apr. 2022.

\bibitem{rebsamen2020cortical}
M.~Rebsamen, C.~Rummel, M.~Reyes, R.~Wiest, and R.~McKinley,
``Direct cortical thickness estimation using deep learning-based anatomy
segmentation and cortex parcellation,''
\emph{Hum. Brain Mapp.}, vol.~41, no.~17, pp.~4804--4814, 2020.

\bibitem{islam2025aibrain}
K.~T. Islam, S.~Zhong, P.~Zakavi, H.~Kavnoudias, S.~Farquharson,
G.~Durbridge, M.~Barth, A.~Dwyer, K.~L. McMahon, P.~M. Parizel,
R.~McIntyre, G.~F. Egan, M.~Law, and Z.~Chen,
``AI improves consistency in regional brain volumes measured in
ultra-low-field MRI and 3T MRI,''
\emph{Front. Neuroimaging}, vol.~4, Art.~no.~1588487, Jun. 2025.

\bibitem{javadi2025diffusion}
M.~Javadi, R.~Griffin, P.~Tsiamyrtzis, E.~Leiss, A.~G. Webb,
and N.~V. Tsekos,
``In-silico comparison of a diffusion model with conventionally trained
deep networks for translating 64\,mT to 3T brain FLAIR,''
\emph{Sci. Rep.}, vol.~15, Art.~no.~21806, 2025.

\bibitem{gopinath2025cortical}
K.~Gopinath, A.~Sorby-Adams, J.~W. Ramirez, D.~Zemlyanker,
J.~Guo, D.~Hunt, C.~L. Mac~Donald, C.~D. Keene, T.~Coalson,
M.~F. Glasser, D.~Van~Essen, M.~S. Rosen, O.~Puonti,
W.~T. Kimberly, and J.~E. Iglesias,
``From low field to high value: Robust cortical mapping from
low-field MRI,''
arXiv preprint arXiv:2505.12228, May 2025.

\bibitem{minocha2024lfmri}
A.~Minocha and S.~Yang,
``Low field MRI deep learning framework for non-hospitalizing early
detection and characterization of Alzheimer's disease pathology,''
in \emph{Proc. IEEE Int. Conf. Big Data (BigData)}, 2024,
pp.~1--8. doi:10.1109/BigData62323.2024.10826144.

\bibitem{dayarathna2024adversarial}
S.~Dayarathna, K.~T. Islam, and Z.~Chen,
``Ultra low-field to high-field MRI translation using adversarial
diffusion,''
in \emph{Proc. IEEE ISBI}, 2024, pp.~1--4.
doi:10.1109/ISBI56570.2024.10635808.

\bibitem{baljer2024paediatric}
L.~Baljer, Y.~Zhang, N.~J. Bourke, K.~A. Donald, L.~E. Bradford,
J.~E. Ringshaw, S.~R. Williams, S.~C.~L. Deoni, S.~C.~R. Williams,
Khula SA Study Team, F.~V{\'a}{\v s}a, and R.~J. Moran,
``Ultra-low-field paediatric MRI in low- and middle-income countries:
Super-resolution using a multi-orientation {U-Net},''
\emph{Hum. Brain Mapp.}, vol.~46, Art.~no.~e70112, 2025.

\end{thebibliography}

\end{document}
